% GENERATED FILE. Edit canonical lessons, then run npm run generate:books.
% canonical-source-hash: fnv1a64:fba2b6434288f41c
% canonical-lessons: ZH-C250-aihao, ZH-C250-xingqu, ZH-C250-xiaoshuo, ZH-C250-gushi, ZH-C250-xiaqi

\chapter{Hobbies}
\label{ch:zh-hobbies}
\hlchaptermodality{250}

\begin{chapteropening}
\textbf{By the end of this chapter:} \emph{I can ask about hobbies and interests, and talk about novels, stories and playing chess.}

Complete the last lesson of chapter 250: I can ask about hobbies and interests, and talk about novels, stories and playing chess.
\end{chapteropening}

\section[àihào]{àihào — hobby}

\label{lesson:ZH-C250-aihao}



\begin{quote}
Before the new one: say the Mandarin for marks at school, then the Mandarin for hard-working.
\end{quote}

\subsection*{You'll want to know: àihào}
\emph{àihào} — \textquotedblleft{}hobby\textquotedblright{}.

\begin{quote}
\emph{nǐ yǒu shénme àihào} — What hobbies do you have?
\end{quote}

Its characters come later; for now it is a word for the ear.

\begin{cousinweb}
\emph{àihào} is built from \emph{ài} 'to love' and \emph{hào} 'to be fond of', both inherited from Middle Chinese, \emph{'ojH} and \emph{xawH} (Guangyun, in Baxter's transcription), which glosses the first 'to cherish' and the second with this very word. The second is \zh{好} said with a falling tone.
\end{cousinweb}

\begin{grammarlens}[title={What you've built}]
One more word for what I like.
\end{grammarlens}

\subsection*{Your turn}
\begin{itemize}
\raggedright
  \item \emph{Say it:} \emph{àihào}
  \item \emph{Say it:} \emph{àihào}, once more
  \item \emph{Say it:} \emph{nǐ yǒu shénme àihào} — What hobbies do you have?
  \item \emph{Recall:} say the Mandarin for marks at school, then the Mandarin for hard-working, then say \emph{àihào} again
\end{itemize}

\begin{tcolorbox}[breakable,colback=teal!4,colframe=teal!35!black,title={Before you move on}]
What does àihào mean? (\textquotedblleft{}Hobby\textquotedblright{}.) Say it once more.
\end{tcolorbox}
\section[xìngqù]{xìngqù — interest; interest in something}

\label{lesson:ZH-C250-xingqu}



\begin{quote}
Before the new one: say the Mandarin for hard-working, then the Mandarin for a hobby.
\end{quote}

\subsection*{You'll want to know: xìngqù}
\emph{xìngqù} — \textquotedblleft{}interest; interest in something\textquotedblright{}.

\begin{quote}
\emph{wǒ méiyǒu xìngqù} — I'm not interested.
\end{quote}

Its characters come later; for now it is a word for the ear.

\begin{cousinweb}
\emph{xìngqù} is built from \emph{xìng} and \emph{qù}, both inherited from Middle Chinese, \emph{xingH} and \emph{tshjuH} (Guangyun, in Baxter's transcription). The dictionary glosses the second 'to tend towards'; it gives no sense for the first in this tone, and its level-tone reading means 'to rise'. How the pair came to mean an interest is not settled.
\end{cousinweb}

\begin{grammarlens}[title={What you've built}]
One more word for what I like.
\end{grammarlens}

\subsection*{Your turn}
\begin{itemize}
\raggedright
  \item \emph{Say it:} \emph{xìngqù}
  \item \emph{Say it:} \emph{xìngqù}, once more
  \item \emph{Say it:} \emph{wǒ méiyǒu xìngqù} — I'm not interested.
  \item \emph{Recall:} say the Mandarin for hard-working, then the Mandarin for a hobby, then say \emph{xìngqù} again
\end{itemize}

\begin{tcolorbox}[breakable,colback=teal!4,colframe=teal!35!black,title={Before you move on}]
What does xìngqù mean? (\textquotedblleft{}Interest; interest in something\textquotedblright{}.) Say it once more.
\end{tcolorbox}
\section[\texorpdfstring{\zh{小说} (xiǎoshuō)}{小说 (xiǎoshuō)}]{\zh{小说} (xiǎoshuō) — novel; story (in a book)}

\label{lesson:ZH-C250-xiaoshuo}



\begin{quote}
Before the new one: say the Mandarin for a hobby, then the Mandarin for an interest.
\end{quote}

\subsection*{You'll want to know: \zh{小说}}
\textbf{\zh{小说}} — \emph{xiǎoshuō} — \textquotedblleft{}novel; story (in a book)\textquotedblright{}.

\begin{quote}
\textbf{\zh{我喜欢读小说}} \emph{wǒ xǐhuan dú xiǎoshuō} — I like reading novels.
\end{quote}

\begin{cousinweb}
\textbf{\zh{小说}} (\emph{xiǎoshuō}) is built from \zh{小} 'small' and \zh{说} 'to tell', both inherited from Middle Chinese, \emph{sjewX} and \emph{sywet} (Guangyun, in Baxter's transcription), which glosses the second 'to tell'. How 'small tellings' came to mean a novel is not settled.
\end{cousinweb}

\begin{grammarlens}[title={What you've built}]
One more word for what I like.
\end{grammarlens}

\subsection*{Your turn}
\begin{itemize}
\raggedright
  \item \emph{Say it:} \emph{xiǎoshuō}
  \item \emph{Say it:} \emph{xiǎoshuō}, once more
  \item \emph{Say it:} \emph{wǒ xǐhuan dú xiǎoshuō} — I like reading novels.
  \item \emph{Recall:} say the Mandarin for a hobby, then the Mandarin for an interest, then say \emph{xiǎoshuō} again
\end{itemize}

\begin{tcolorbox}[breakable,colback=teal!4,colframe=teal!35!black,title={Before you move on}]
What does \zh{小说} mean? (\textquotedblleft{}Novel; story (in a book)\textquotedblright{}.) Say it once more.
\end{tcolorbox}
\section[gùshi]{gùshi — story}

\label{lesson:ZH-C250-gushi}



\begin{quote}
Before the new one: say the Mandarin for an interest, then the Mandarin for a novel.
\end{quote}

\subsection*{You'll want to know: gùshi}
\emph{gùshi} — \textquotedblleft{}story\textquotedblright{}.

\begin{quote}
\emph{zhège gùshi hěn yǒu yìsi} — This story is really interesting.
\end{quote}

Its characters come later; for now it is a word for the ear.

\begin{cousinweb}
\emph{gùshi} is built from \emph{gù} 'old' and \zh{事} 'an affair', both inherited from Middle Chinese, \emph{kuH} and \emph{dzriH} (Guangyun, in Baxter's transcription), which glosses the first 'old'. In the word the second loses its tone.
\end{cousinweb}

\begin{grammarlens}[title={What you've built}]
One more word for what I like.
\end{grammarlens}

\subsection*{Your turn}
\begin{itemize}
\raggedright
  \item \emph{Say it:} \emph{gùshi}
  \item \emph{Say it:} \emph{gùshi}, once more
  \item \emph{Say it:} \emph{zhège gùshi hěn yǒu yìsi} — This story is really interesting.
  \item \emph{Recall:} say the Mandarin for an interest, then the Mandarin for a novel, then say \emph{gùshi} again
\end{itemize}

\begin{tcolorbox}[breakable,colback=teal!4,colframe=teal!35!black,title={Before you move on}]
What does gùshi mean? (\textquotedblleft{}Story\textquotedblright{}.) Say it once more.
\end{tcolorbox}
\section[xiàqí]{xiàqí — to play chess}

\label{lesson:ZH-C250-xiaqi}



\begin{quote}
Before the new one: say the Mandarin for a novel, then the Mandarin for a story.
\end{quote}

\subsection*{You'll want to know: xiàqí}
\emph{xiàqí} — \textquotedblleft{}to play chess\textquotedblright{}.

\begin{quote}
\emph{wǒmen xiàwǔ yìqǐ xiàqí} — We're playing chess together this afternoon.
\end{quote}

Its characters come later; for now it is a word for the ear.

\begin{cousinweb}
\emph{xiàqí} is built from \zh{下} 'to go down' and \emph{qí} 'a board game', both inherited from Middle Chinese, \emph{haeX} and \emph{gi} (Guangyun, in Baxter's transcription), which glosses the first 'to descend' and lists the second under another spelling of the character, quoting an older book on who first made the board game weiqi.
\end{cousinweb}

\begin{grammarlens}[title={What you've built}]
One more word for what I like.
\end{grammarlens}

\subsection*{Your turn}
\begin{itemize}
\raggedright
  \item \emph{Say it:} \emph{xiàqí}
  \item \emph{Say it:} \emph{xiàqí}, once more
  \item \emph{Say it:} \emph{wǒmen xiàwǔ yìqǐ xiàqí} — We're playing chess together this afternoon.
  \item \emph{Recall:} say the Mandarin for a novel, then the Mandarin for a story, then say \emph{xiàqí} again
\end{itemize}

\begin{tcolorbox}[breakable,colback=teal!4,colframe=teal!35!black,title={Before you move on}]
What does xiàqí mean? (\textquotedblleft{}To play chess\textquotedblright{}.) Say it once more.
\end{tcolorbox}
